\subsection{QUOTIENTS OF EQUIVALENCE RELATIONS}

Let R, S be two equivalence relations with respect to two letters x, y. We shall say that S is finer than R (or that R is coarser than S) if the relation \(S \Longrightarrow R\) is true. If R and S are equivalence relations on the same set E, the statement that S is finer than R means that the graph of S is contained in that of R, or again that every equivalence class with respect to S is contained in an equivalence class with respect to R; or, equivalently, that every equivalence class with respect to R is saturated with respect to S.

Examples

\begin{enumerate}
\def\labelenumi{(\arabic{enumi})}
\tightlist
\item
  The relation `` \(x \in E\) and \(y \in E\) and \(x = y\) '' is finer than every equivalence relation on E. The relation `` \(x \in E\) and \(y \in E\) '' is coarser than every equivalence relation on E.
\end{enumerate}

* (2) The equivalence relation `` \(x \in \mathbf{Z}\) and \(y \in \mathbf{Z}\) and \(x - y\) is divisible by 4'' is finer than the equivalence relation `` \(x \in \mathbf{Z}\) and \(y \in \mathbf{Z}\) and \(x - y\) is divisible by 2''. *

Let R and S be two equivalence relations on the same set E, such that S is finer than R. Let f and g be the canonical mappings of E onto E/R and E onto E/S respectively; then the function f is compatible with S. Let h be the function induced by f on passing to the quotient with respect to S; then h is a mapping of E/S onto E/R. The equivalence relation associated with h on E/S is called the quotient of R by S and is denoted by R/S. The relation \(x \equiv y \pmod{R}\) is equivalent to \(g(x) \equiv g(y) \pmod{R/S}\) , and the equivalence classes with respect to R/S are the images under g of the equivalence classes with respect to R. Let \(h = j \circ h_{2} \circ h_{1}\) be the canonical decomposition (no. 5) of the mapping h. Then \(h_{1}\) is the canonical mapping of E/S onto (E/S)/(R/S), j is the identity mapping of E/R, and \(h_{2}\) is a one-to-one mapping of (E/S)/(R/S) onto E/R. The mapping \(h_{2}\) and its inverse are said to be canonical.

Conversely, consider an arbitrary equivalence relation T on the set E/S, and let R be the equivalence relation on E which is the inverse image under g of the relation T (no. 6). Since the relation \(x \equiv y \pmod{R}\) is equivalent to \(g(x) \equiv g(y) \pmod{T}\) , it follows that T is equivalent to R/S.

